\begin{abstract}
Standard analyses of catastrophic health expenditure (CHE) in developing countries rely on binary threshold indicators that mask heterogeneity across the spending distribution. Using nationally representative data from Peru's 2024 National Household Survey (ENAHO, $N = 33{,}691$ households), we apply conditional quantile regression at $\tau = \{0.50, 0.75, 0.90\}$ to examine whether the correlates of out-of-pocket (OOP) health expenditure share vary across the conditional distribution. We interact SIS (Seguro Integral de Salud) membership with consumption quintiles, restricting the primary analysis to the SIS--uninsured comparison given data limitations in EsSalud coverage measurement. Our results reveal a monotonic regressive protection gradient: at the median, SIS membership shows no statistically significant association with OOP share, but at upper quantiles ($\tau = 0.75, 0.90$), SIS--quintile interactions indicate larger OOP reductions for higher-quintile enrollees---a pattern consistent with regressive financial protection rather than a middle-income squeeze. Consumption quintile gradients steepen from $\tau = 0.50$ to $\tau = 0.90$, though this gradient is partly mechanical given that quintiles are constructed from total consumption inclusive of health spending. These results are provisional: bootstrap inference uses 200 replications (below the 999 standard), and the 44.1\% zero-OOP mass limits the range of informative quantiles. Subject to these caveats, the findings suggest that mean-based CHE analyses understate the concentration of financial risk in the upper tail of the spending distribution.

\medskip
\noindent \textbf{Keywords:} catastrophic health expenditure, quantile regression, health insurance, out-of-pocket spending, Peru, ENAHO

\noindent \textbf{JEL codes:} I13, I14, I18, C21
\end{abstract}
